\documentclass[11pt, a4paper]{amsart}
\usepackage{amsmath, amssymb, amsthm, amsfonts}
\usepackage{mathtools}
\usepackage[margin=1in]{geometry}
\usepackage{hyperref}

\newtheorem{theorem}{Theorem}[section]
\newtheorem{lemma}[theorem]{Lemma}
\newtheorem{proposition}[theorem]{Proposition}
\newtheorem{corollary}[theorem]{Corollary}
\newtheorem{fact}[theorem]{Standard Fact}
\theoremstyle{definition}
\newtheorem{definition}[theorem]{Definition}
\newtheorem{remark}[theorem]{Remark}
\newtheorem{notation}[theorem]{Notation}
\newtheorem{hypothesis}[theorem]{Hypothesis}
\newtheorem{axiom}[theorem]{Axiom}

\newcommand{\R}{\mathbb{R}}
\newcommand{\C}{\mathbb{C}}
\newcommand{\Q}{\mathbb{Q}}
\newcommand{\Z}{\mathbb{Z}}
\newcommand{\HH}{\mathbb{H}}
\newcommand{\Irr}{\operatorname{Irr}}
\newcommand{\rk}{\operatorname{rank}}
\newcommand{\Hmin}{\mathbf{H}_{\min}}
\newcommand{\Hor}{\mathbf{H}_{\mathrm{or}}}
\newcommand{\Hdist}{\mathbf{H}_{\mathrm{dist}}}
\newcommand{\Hquot}{\mathbf{H}_{\mathrm{quot}}}
\newcommand{\Hfifty}{\mathbf{H}_{50}}
\newcommand{\Hgauge}{\mathbf{H}_{\mathrm{gauge}}}
\newcommand{\Ophi}{\mathbb{Z}[\varphi]}
\newcommand{\Wph}{W_{\mathrm{ph}}}
\newcommand{\Nslot}{N_{\mathrm{slot}}}
\newcommand{\kappaTwoI}{\kappa_{2I}}
\newcommand{\Phiraw}{\Phi_{\mathrm{raw}}}
\newcommand{\PhiHfour}{\Phi_{H_4}}
\newcommand{\Irrrat}{\Irr_{\mathrm{rat}}}

\title{The Fine-Structure Theorem:\\
  $\alpha^{-1} = 137 + \pi/87$\\
  as a Conditional Identity in the $\alpha$-Chain}
\author{}
\date{}

\begin{document}
\maketitle

\begin{abstract}
This is the third and final entry of the $\alpha$-chain in
the cascade-correspondence programme. It assembles the
inherited results of $\alpha$-1 (\emph{The 12D Closure
Theorem}) and $\alpha$-2 (\emph{Phason--Coxeter Alignment})
together with classical icosian / 600-cell / Coxeter theory
into the conditional mathematical identity
\[
  \alpha^{-1} \;=\; 137 \;+\; \pi/87,
\]
where the integer $137 = 87 + 50$ decomposes as
\begin{itemize}
\item $87 = \Nslot$, the phason slot count from $\alpha$-2's
  Theorem $T_B$ (slot-count clause), citable conditionally on
  $\alpha$-2's hypothesis stack
  $\Hmin + \Hor + \Hdist + \Hquot$;
\item $50 = \kappaTwoI$, the sum of nontrivial rational-sector
  Laplacian eigenvalues of the binary icosahedral group $2I$
  acting on the $600$-cell graph, identified with $50$ under
  a named structural hypothesis $\Hfifty$
  (Hypothesis~\ref{hyp:H50}).
\end{itemize}
The phase term $\pi/87$ is derived from the per-slot
distribution of the $H_4$-observable phase
$\PhiHfour = 88\pi$ across $\Nslot = 87$ slots, after
separating a center-gauge contribution of magnitude $\pi$.
The center-gauge separation is encoded as a named structural
hypothesis $\Hgauge$ (Hypothesis~\ref{hyp:Hgauge}).

The final identity (Theorem~\ref{thm:Talpha}) is conditional
on the full seven-hypothesis stack
\[
  \Hmin \,+\, \Hor \,+\, \Hdist \,+\, \Hquot \,+\,
  \Hfifty \,+\, \Hphasehalve \,+\, \Hgauge,
\]
where $\Hphasehalve$ (Hypothesis~\ref{hyp:phase-halving})
asserts that the classical $E_8 \to H_4$ root-layer $2$-to-$1$
projection transfers to a factor-of-$2$ halving of the
cumulative Coxeter phase.
We make no claim that this identity is unconditional, that
$\alpha^{-1}$ here equals an experimental observable, that
this paper proves any law of nature, or that the
hypothesis stack is itself reducible to lower-level
mathematics within $\alpha$-3. The paper's deliverable is
the strictly mathematical conditional identity above.

We make no claim about photons, microtubules, biology, the
Standard Model, or any phenomenology. We do not depend on the
infrastructure papers P1--P7. We do not cite the
\texttt{cascade-derivation/} working notes,
\texttt{paper-xxii}, \texttt{paper-xxxiv}, or any other
downstream VFD manuscript as proof source.
\end{abstract}

\tableofcontents

%=====================================================================
\section{Introduction}
\label{sec:intro}
%=====================================================================

\subsection{Position of the paper}

This is the third and final entry of the $\alpha$-chain in
the cascade-correspondence programme: a strictly mathematical
synthesis paper that combines the inherited results of
$\alpha$-1 \cite{Cascade12D} and $\alpha$-2
\cite{CascadePhasonCoxeter} with classical icosian /
$600$-cell theory into a conditional identity for $\alpha^{-1}$
as a \emph{derived mathematical constant}.

The paper is parallel to the infrastructure stack P1--P7 and
does not depend on it. It does not depend on any
\texttt{cascade-derivation/} working note, on any downstream
VFD paper, or on any computer-algebra session.

\subsection{What this paper proves}

\textbf{Inherited (cited from $\alpha$-1, $\alpha$-2).}
\begin{description}
\item[\textbf{$\alpha$-1, $C_1$, $C_2$, $C_3$}
  \cite{Cascade12D}.] $C_2$ unconditional; $C_1$, $C_3$
  conditional on $\Hmin$.
\item[\textbf{$\alpha$-2, $T_A$, $T_B$-$88$, $T_B$-$87$}
  \cite{CascadePhasonCoxeter}.] $T_A$ and $T_B$-$88$
  conditional on $\Hmin + \Hor + \Hdist$; $T_B$-$87$
  conditional on $\Hmin + \Hor + \Hdist + \Hquot$.
\end{description}

\textbf{Classical inputs.}
\begin{description}
\item[\textbf{Lemma~\ref{lem:McKay-recap}.}] McKay
  correspondence $2I \leftrightarrow \widetilde{E}_8$
  (recap from $\alpha$-2).
\item[\textbf{Lemma~\ref{lem:c15}.}] $c^{15} = -I$ in
  $W(E_8)$, the standard $\Z/2$ central involution generated
  by powers of a Coxeter element.
\item[\textbf{Lemma~\ref{lem:E8-to-H4-2to1}.}] The classical
  $E_8 \to H_4$ cut-and-project package is two-to-one on the
  minimal-vector / root layer (Moody--Patera). Transfer of
  this root-layer datum to a halving of the cumulative phase
  bookkeeping is not classical; it is the content of
  Hypothesis~\ref{hyp:phase-halving} ($\Hphasehalve$), and is
  used only under that hypothesis.
\end{description}

\textbf{Layer A (definitions inheriting the full $\alpha$-1
+ $\alpha$-2 stack $\Hmin + \Hor + \Hdist + \Hquot$).}
\begin{description}
\item[\textbf{Definition~\ref{def:Nslot}.}] $\Nslot := 87$
  cited from $\alpha$-2's $T_B$-$87$
  (Fact~\ref{fact:imports}(e), under
  $\Hmin + \Hor + \Hdist + \Hquot$).
\item[\textbf{Definitions~\ref{def:phi-raw}, \ref{def:phi-H4}.}]
  Total Coxeter phase $\Phiraw = 88 \cdot 2\pi = 176\pi$
  (from $\alpha$-2's $T_B$-$88$, conditional on
  $\Hmin + \Hor + \Hdist$); two-to-one $H_4$-observable
  phase $\PhiHfour = \Phiraw / 2 = 88\pi$ (using
  Lemma~\ref{lem:E8-to-H4-2to1} \emph{together with}
  Hypothesis~\ref{hyp:phase-halving}, which transfers the
  root-layer 2-to-1 to the phase bookkeeping).
\end{description}

\textbf{Layer B (conditional on $\alpha$-2 stack + $\Hfifty$).}
\begin{description}
\item[\textbf{Theorem~\ref{thm:T137}}] (\textbf{$T_{137}$},
  integer decomposition): $137 = \Nslot + \kappaTwoI = 87
  + 50$, conditional on $\Hmin + \Hor + \Hdist + \Hquot +
  \Hfifty$.
\end{description}

\textbf{Layer C (conditional on $\alpha$-2 stack + $\Hphasehalve + \Hgauge$).}
\begin{description}
\item[\textbf{Theorem~\ref{thm:Tpi}}] (\textbf{$T_{\pi/87}$},
  per-slot phase): $\PhiHfour / \Nslot = \pi + \pi/87$,
  with the $\PhiHfour = 88\pi$ identification requiring
  $\Hphasehalve$ and the $+\pi$ identification as the $\Z/2$
  center gauge under $\Hgauge$. Conditional on
  $\Hmin + \Hor + \Hdist + \Hquot + \Hphasehalve + \Hgauge$.
\end{description}

\textbf{Layer D (final synthesis, conditional on the full
seven-hypothesis stack).}
\begin{description}
\item[\textbf{Theorem~\ref{thm:Talpha}}] (\textbf{$T_\alpha$}):
  $\alpha^{-1} = 137 + \pi/87$ as a derived mathematical
  constant. Conditional on $\Hmin + \Hor + \Hdist + \Hquot
  + \Hfifty + \Hphasehalve + \Hgauge$.
\end{description}

\subsection{What this paper does NOT claim}

\begin{enumerate}
\item No claim that $T_\alpha$ is unconditional. Every
  hypothesis in the stack is named and tracked.
\item No claim that $\alpha^{-1}$ as defined here equals the
  experimental fine-structure constant. Inside the proof
  layer, $\alpha^{-1}$ is a mathematical quantity derived from
  $\Nslot$, $\kappaTwoI$, and the residual phase term. Any
  comparison with empirical measurement is outside the proof
  layer.
\item No claim about photons, microtubules, biology, the
  Standard Model, or any phenomenology.
\item No claim that $\Hfifty$ or $\Hgauge$ is itself a
  theorem of $\alpha$-3. Both are named structural
  hypotheses with explicit motivation
  (Remarks~\ref{rmk:H50-status}, \ref{rmk:Hgauge-status});
  reducing them to upstream mathematics is deferred.
\item No first-principles-physics rhetoric.
\item No dependence on the infrastructure papers P1--P7, on
  \texttt{cascade-derivation/} working notes, on
  \texttt{paper-xxii}, \texttt{paper-xxxiv}, or on any
  downstream VFD manuscript.
\item No computer-algebra session is cited as proof source;
  any GAP / numerical cross-check appears only in non-load-bearing
  remarks.
\end{enumerate}

\subsection{Source audit}

\textbf{Internal.}
\begin{itemize}
\item Axiom $\mathbf{F1}$ from
  \cite[\S F1]{CascadeFoundations}: chain provenance only.
\item $\alpha$-1 \cite{Cascade12D}: $C_2$ unconditional;
  $C_1$, $C_3$ conditional on $\Hmin$.
\item $\alpha$-2 \cite{CascadePhasonCoxeter}: $T_A$,
  $T_B$-$88$ conditional on $\Hmin + \Hor + \Hdist$;
  $T_B$-$87$ conditional on $\Hmin + \Hor + \Hdist + \Hquot$.
\end{itemize}

\textbf{Classical.}
\begin{itemize}
\item McKay 1980 \cite{McKay80}: ADE / $2I$ /
  $\widetilde{E}_8$ correspondence.
\item Conway--Sloane \cite{ConwaySloane} Ch.~8: icosian
  ring, $E_8$ embedding, $H_4$ / $600$-cell interface.
\item Moody--Patera \cite{MoodyPatera}: $E_8 \to H_4$
  cut-and-project.
\item Bourbaki \cite{Bourbaki46} Ch.~4--6 / Humphreys
  \cite{Humphreys} \S 3.16--3.20: Coxeter element
  eigenvalues, exponents, $E_8$ data.
\item Du Val \cite{DuVal} Ch.~3: standard $2I$ character
  table.
\item In-paper Proposition~\ref{prop:class-sum}: class-sum /
  Cayley-graph spectral formula (proved here for self-contained
  reference).
\end{itemize}

%=====================================================================
\section{Preliminaries: Imports from $\alpha$-1, $\alpha$-2,
  and Classical Sources}
\label{sec:prelims}
%=====================================================================

\subsection{Inherited convention}

We adopt $\alpha$-2's resolved convention
\cite[Notation~not:rank-conv]{CascadePhasonCoxeter}:
\emph{rank} means $\rk_\Ophi$ unless explicitly marked
$\rk_\Z$. We do not re-litigate the $\rk_\Z$-vs-$\rk_\Ophi$
distinction; for $\alpha$-3 the load-bearing inherited
quantities are integers and real numbers
($\Nslot$, $\Wph$, etc.), so the convention does not enter
visibly here.

\subsection{Inherited theorem statements}

\begin{fact}[Imports from $\alpha$-1 and $\alpha$-2]
\label{fact:imports}
We freely cite the following.
\begin{enumerate}
\item[(a)] $\alpha$-1 Theorem $C_2$ (\emph{unconditional}):
  the upward cascade closes at $\Z$-rank $12$
  \cite[Thm.~thm:C2]{Cascade12D}.
\item[(b)] $\alpha$-1 Hypothesis $\Hmin$
  (\emph{conditional}): the phason coefficient ring is
  $\Ophi$ \cite[Hyp.~hyp:Hmin]{Cascade12D}.
\item[(b1)] $\alpha$-1 Theorem $C_1$ (\emph{conditional on
  $\Hmin$}): the phason complement satisfies $M \cong
  \Ophi^{2}$ as $\Ophi$-module
  \cite[Thm.~thm:C1]{Cascade12D}.
\item[(b2)] $\alpha$-1 Theorem $C_3$ (\emph{conditional on
  $\Hmin$}): the phason coefficient ring at $L_{12}$
  closure equals $\Ophi$
  \cite[Thm.~thm:C3]{Cascade12D}.
\item[(c)] $\alpha$-2 Theorem $T_A$ (\emph{conditional on
  $\Hmin + \Hor + \Hdist$, relative to a fixed Coxeter
  element $c \in W(E_8)$}): the Coxeter-resolved phason
  module $M_{\rm Cox}(c) \cong \Ophi^{4}$ has one
  $\Ophi$-line per Coxeter plane $P_k(c)$
  \cite[Thm.~thm:TA]{CascadePhasonCoxeter}.
\item[(d)] $\alpha$-2 Theorem $T_B$-$88$ (\emph{conditional
  on $\Hmin + \Hor + \Hdist$}): the total phason winding is
  $\Wph = \sum_{k=1}^{4} m_k^+ = 17 + 19 + 23 + 29 = 88$
  \cite[Thm.~thm:TB-88]{CascadePhasonCoxeter}.
\item[(e)] $\alpha$-2 Theorem $T_B$-$87$ (\emph{conditional
  on $\Hmin + \Hor + \Hdist + \Hquot$}): the phason slot
  count is $\Nslot = \Wph - 1 = 87$
  \cite[Thm.~thm:TB-87]{CascadePhasonCoxeter}.
\end{enumerate}
\end{fact}

\begin{definition}[Phason slot count, inherited]
\label{def:Nslot}
$\Nslot := 87$, cited from $\alpha$-2 Theorem $T_B$-$87$
(Fact~\ref{fact:imports}(e)). The integer $\Nslot$ is used
in $\alpha$-3 only as an inherited quantity; it carries the
$\alpha$-2 conditional fence
$\Hmin + \Hor + \Hdist + \Hquot$.
\end{definition}

%=====================================================================
\section{The $2I$ / $600$-cell / McKay Package}
\label{sec:McKay-pack}
%=====================================================================

\begin{lemma}[McKay correspondence: $2I \leftrightarrow
  \widetilde{E}_8$, recap]
\label{lem:McKay-recap}
The binary icosahedral group $2I \subset SU(2)$ corresponds
to the affine Dynkin diagram $\widetilde{E}_8$ under the
McKay correspondence: the irreducible representations of
$2I$ are in canonical bijection with the nodes of
$\widetilde{E}_8$ (including the trivial representation on
the affine node).
\end{lemma}

\begin{proof}
Classical McKay 1980 \cite[\S 3, pp.~183--186]{McKay80};
recap of \cite[Lem.~lem:McKay]{CascadePhasonCoxeter}.
\end{proof}

\begin{definition}[$2I$ rational sector $\Irrrat(2I)$]
\label{def:Irr-rat}
We define $\Irrrat(2I) \subset \Irr(2I)$ as the set of
\emph{complex irreducibles of $2I$ with rational-valued
character}, i.e., $\rho \in \Irr(2I)$ such that
$\chi_\rho(g) \in \Q$ for all $g \in 2I$. The nontrivial
rational sector is $\Irrrat^\times(2I) := \Irrrat(2I)
\setminus \{\mathbf{1}\}$ (the trivial representation
$\mathbf{1}$ is excluded by definition).
\end{definition}

\begin{remark}[Why ``rational''?]
\label{rmk:Irr-rat-justif}
The choice of rational sector is the algebraic restriction
under which the Laplacian spectrum (Def.~\ref{def:Lambda})
gives integer-valued plus rational-valued eigenvalues, which
will sum to an integer $\kappaTwoI$. Other restrictions of
the spectrum (full complex sector, real sector) give
non-integer sums and do not align with the integer $50$
target. The choice ``rational sector'' is a definitional
restriction, not derived; we make it explicit so that
$\Hfifty$ below has unambiguous content.
\end{remark}

\begin{definition}[$600$-cell graph and Laplacian on a
  $2I$-isotypic component]
\label{def:Lambda}
Let $\Gamma$ be the $1$-skeleton of the $600$-cell viewed as
a $12$-regular Cayley graph of $2I$ with respect to a fixed
neighbour conjugacy class $S \subset 2I$ of size $|S| = 12$
(the $12$ icosahedral neighbours). Let $L = 12 I - A$
be the unnormalised graph Laplacian, where $A$ is the
adjacency matrix of $\Gamma$.

For each complex irreducible $\rho \in \Irr(2I)$, the
isotypic component $V_\rho \subset \C[2I]$ is
$L$-invariant, and $L|_{V_\rho}$ acts as a scalar
$\lambda_\rho \in \R$ given by the class-sum formula stated
in Proposition~\ref{prop:class-sum} below.
\end{definition}

\begin{proposition}[Class-sum / Cayley-graph spectral
  formula]
\label{prop:class-sum}
Let $G$ be a finite group, $S \subset G$ a symmetric
($S = S^{-1}$) union of conjugacy classes of $G$, and
$\Gamma = \mathrm{Cay}(G, S)$ the Cayley graph of $G$ with
respect to $S$. Let $A$ be the adjacency matrix and $L =
|S| \cdot I - A$ the unnormalised Laplacian. Decompose the
regular representation as
\[
  \C[G] \;\cong\; \bigoplus_{\rho \in \Irr(G)}
    W_\rho^{\oplus \dim \rho},
\]
where $W_\rho$ is the irreducible carrier of $\rho$ (with
$\dim W_\rho = \dim \rho$), so that the $\rho$-\emph{isotypic
component} of $\C[G]$ is $V_\rho := W_\rho^{\oplus \dim \rho}$
(of total $\C$-dimension $(\dim \rho)^2$).
Then on each irreducible carrier $W_\rho$ (and hence as the
same scalar on every copy in $V_\rho$), $A$ and $L$ act as
scalars $a_\rho, \lambda_\rho$ given by
\[
  a_\rho \;=\; \frac{\sum_{s \in S} \chi_\rho(s)}{\dim \rho},
  \qquad
  \lambda_\rho \;=\; |S| - a_\rho.
\]
When $S$ is a single conjugacy class with representative
$s_0 \in S$, this simplifies to $a_\rho = |S| \cdot
\chi_\rho(s_0) / \dim \rho$, hence
\[
  \lambda_\rho \;=\; |S| \cdot \left(1 -
    \frac{\chi_\rho(s_0)}{\dim \rho}\right).
\]
\end{proposition}

\begin{proof}
The class sum $\Sigma := \sum_{s \in S} s$ lies in the centre
of $\C[G]$, hence by Schur's lemma acts on the irreducible
carrier $W_\rho$ as a scalar $a_\rho \in \C$. Computing the
trace of $\Sigma$ on $W_\rho$:
\[
  \mathrm{tr}_{W_\rho}\bigl(\rho(\Sigma)\bigr)
  \;=\; \sum_{s \in S} \chi_\rho(s),
\]
while the trace of the scalar $a_\rho \cdot I_{W_\rho}$ on
$W_\rho$ is $a_\rho \cdot \dim W_\rho = a_\rho \cdot \dim
\rho$. Hence
\[
  a_\rho \;=\;
  \frac{\sum_{s \in S} \chi_\rho(s)}{\dim \rho}.
\]
On the isotypic component $V_\rho = W_\rho^{\oplus \dim
\rho}$, $\Sigma$ acts as the same scalar $a_\rho$ on each
copy of $W_\rho$ (since the class sum is central).
The Cayley adjacency matrix is $A = \rho_{\rm reg}(\Sigma)$
(where $\rho_{\rm reg}$ is the regular representation),
so $A$ inherits this scalar action on each $V_\rho$.
The Laplacian formula $\lambda_\rho = |S| - a_\rho$ is
direct. The single-class simplification uses that $\chi_\rho$
is constant on conjugacy classes.
\end{proof}

\begin{remark}[Application to $\Gamma$ and the $600$-cell]
\label{rmk:Lambda-application}
For $G = 2I$ and $S = $ the $12$-element neighbour conjugacy
class on the $600$-cell with representative $s_0 \in S$,
Proposition~\ref{prop:class-sum} gives
\[
  \lambda_\rho \;=\; 12 - \frac{12 \, \chi_\rho(s_0)}{\dim
    \rho}.
\]
This is the formula entering Def.~\ref{def:Lambda} and the
sum $\kappaTwoI$ of Def.~\ref{def:kappa}.
\end{remark}

\begin{definition}[Sum $\kappaTwoI$]
\label{def:kappa}
The \emph{rational-sector Laplacian sum} is
\[
  \kappaTwoI \;:=\; \sum_{\rho \in \Irrrat^\times(2I)}
    \lambda_\rho \;\in\; \R.
\]
\end{definition}

\begin{lemma}[$2I$ rational sector and Laplacian spectrum,
  classical content]
\label{lem:2I-laplacian}
The complex irreducible representations of $2I$ have
dimensions $\{1, 2, 2, 3, 3, 4, 4, 5, 6\}$, with the unique
$1$-dimensional being the trivial representation $\mathbf{1}$
\cite[Ch.~3]{DuVal}. The character table is classical and
records that the rational-character sector
$\Irrrat(2I)$ consists of those irreducibles whose character
values lie in $\Q$ for every conjugacy class. The
Laplacian eigenvalues $\lambda_\rho$ of
Def.~\ref{def:Lambda} are the standard rationals (or
quadratic irrationals over $\Q(\sqrt 5)$ on certain
non-rational classes) given by the class-sum formula.
\end{lemma}

\begin{proof}
Du Val \cite[Ch.~3]{DuVal} for the $2I$ character table; the
class-sum / Cayley-graph spectrum formula is the in-paper
Proposition~\ref{prop:class-sum}. The dimension list and the
rational-character classification follow by direct inspection
of the character table.
\end{proof}

%=====================================================================
\section{Hypothesis $\Hfifty$: the $\kappaTwoI = 50$
  identification}
\label{sec:H50}
%=====================================================================

\begin{hypothesis}[$\Hfifty$: $\kappaTwoI = 50$ with explicit eigenvalue breakdown]
\label{hyp:H50}
The rational-sector Laplacian sum $\kappaTwoI$ of
Def.~\ref{def:kappa}, computed for the $2I$-Cayley structure on
the $600$-cell of Def.~\ref{def:Lambda}, is given by the sum
\[
  \kappaTwoI \;=\; \sum_{\rho \in \Irrrat^\times(2I)}
    \lambda_\rho \;=\; 9 + 12 + 14 + 15 \;=\; 50,
\]
where the four nontrivial rational-sector eigenvalues $9, 12,
14, 15$ are identified with the rational-sector irreps of $2I$
(of dimensions $1, 3, 4, 5$ respectively) via the class-sum
formula $\lambda_\rho = 12 - 12\chi_\rho(s)/\dim\rho$ for the
neighbour class $S$ on the $600$-cell. The hypothesis asserts
both the identification of each individual eigenvalue with a
specific rational-sector irrep and the summation identity.
\end{hypothesis}

\begin{remark}[Status of $\Hfifty$]
\label{rmk:H50-status}
$\Hfifty$ is a \emph{named structural hypothesis}, not a
theorem of $\alpha$-3. To upgrade $\Hfifty$ to a theorem one
would (i) compute the rational-character sector
$\Irrrat^\times(2I)$ explicitly using
Lemma~\ref{lem:2I-laplacian}'s character table, (ii) compute
each $\lambda_\rho$ via the class-sum formula
$\lambda_\rho = 12 - 12\,\chi_\rho(s)/\dim\rho$ for the
neighbour class $S$ on the $600$-cell, and (iii) sum the
results. This computation is in principle finite and
classical; we do not carry it out at full theorem grade in
$\alpha$-3 because doing so requires committing to a single
classical convention for the $2I$ character table, the
$600$-cell adjacency, and the rational-sector definition,
all of which are mildly choice-dependent in the literature
(\cite{DuVal} vs other quaternionic-icosahedral
references). Pinning the conventions to an exact fixed
point and verifying $\kappaTwoI = 50$ thereunder is left
as a cross-check task for a downstream paper. In the
meantime, $\Hfifty$ stays an explicit hypothesis whose
content is exactly the equation $\kappaTwoI = 50$.

A non-load-bearing remark: independent computer-algebra
sessions reportedly give $\kappaTwoI = 50$ under the
conventions specified above; this aligns with $\Hfifty$ but
is not a proof.
\end{remark}

\begin{theorem}[$\mathbf{T_{137}}$: integer decomposition
  $137 = \Nslot + \kappaTwoI$;
  conditional on $\Hmin + \Hor + \Hdist + \Hquot + \Hfifty$]
\label{thm:T137}
Assume the full $\alpha$-2 hypothesis stack
$\Hmin + \Hor + \Hdist + \Hquot$ together with $\Hfifty$.
Then
\[
  \Nslot + \kappaTwoI \;=\; 87 + 50 \;=\; 137.
\]
\end{theorem}

\begin{proof}
By Def.~\ref{def:Nslot} (Fact~\ref{fact:imports}(e),
$\alpha$-2's $T_B$-$87$ under the $\alpha$-2 stack),
$\Nslot = 87$. By $\Hfifty$, $\kappaTwoI = 50$. The
arithmetic $87 + 50 = 137$ is direct.
\end{proof}

\begin{remark}[On $T_{137}$]
\label{rmk:T137-status}
Theorem~\ref{thm:T137} is an \emph{arithmetic} statement:
once both $\Nslot = 87$ (inherited) and $\kappaTwoI = 50$
(hypothesis $\Hfifty$) hold, the sum is $137$ tautologically.
The substantive content of the theorem is therefore the
hypothesis stack: nothing in $\alpha$-3 derives $137$ from
deeper structure; $\alpha$-3 only asserts that $137$ admits
the stated decomposition under the listed hypotheses.
\end{remark}

%=====================================================================
\section{Phase Machinery: Total, Observable, Center Gauge}
\label{sec:phase}
%=====================================================================

\begin{definition}[Total raw phase per Coxeter cycle]
\label{def:phi-raw}
The \emph{total raw phase per Coxeter cycle} is
\[
  \Phiraw \;:=\; \Wph \cdot 2\pi \;=\; 88 \cdot 2\pi
  \;=\; 176\pi,
\]
using $\Wph = 88$ from $\alpha$-2's $T_B$-$88$
(Fact~\ref{fact:imports}(d)).
\end{definition}

\begin{lemma}[$E_8$ minimal-vector layer projects two-to-one
  onto the $H_4$ root system]
\label{lem:E8-to-H4-2to1}
The classical Elser--Sloane cut-and-project from $E_8$ to
$H_4$ projects the $240$ minimal vectors of $E_8$ onto the
$120$ roots of the $H_4$ root system, with each $H_4$ root
arising from a pair $\{v, -v\}$ of $E_8$ minimal vectors
(equivalently, from a $\sigma$-twist-related pair of
icosian / co-icosian points). The factor of $2$ at this
\emph{minimal-vector layer} is the only $2$-to-$1$ datum
needed in this paper to halve the raw Coxeter phase from
$\Phiraw = 176\pi$ to $\PhiHfour = 88\pi$.
\end{lemma}

\begin{proof}
Moody--Patera \cite[Thm.~4.1, Prop.~4.2]{MoodyPatera}
establishes that the cut-and-project from the $E_8$ root
system onto the physical $H_4$ root system maps the $240$
$E_8$-roots onto the $120$ $H_4$-roots with each $H_4$ root
having two preimages related by sign reversal (equivalently,
the $\sigma$-twist on the icosian / co-icosian splitting of
$E_8 = I \oplus \sigma(I)$
\cite[Ch.~8]{ConwaySloane}).
\end{proof}

\begin{remark}[Scope of Lemma~\ref{lem:E8-to-H4-2to1}]
\label{rmk:2to1-scope}
We do \emph{not} claim a global $2$-to-$1$ structure on the
full $H_4$-quasicrystal $Q_{H_4}$ of $\alpha$-1's classical
package; that quasicrystal is a model set
\cite[Def.~def:cap]{Cascade12D}, not an orbit space. The
load-bearing $2$-to-$1$ datum of
Lemma~\ref{lem:E8-to-H4-2to1} is restricted to the
minimal-vector / root layer: $240$ $E_8$-roots project to $120$
$H_4$-roots with sign-paired preimages. This root-layer
statement does \emph{not} automatically transfer to a halving
of phase bookkeeping, which is a different datum; that
transfer is the content of the phase-halving hypothesis
Hypothesis~\ref{hyp:phase-halving} below.
\end{remark}

\begin{hypothesis}[$\Hphasehalve$: root-layer $2$-to-$1$ transfers to the $H_4$-observable phase]
\label{hyp:phase-halving}
The $2$-to-$1$ sign-paired projection from the $E_8$
minimal-vector layer onto the $H_4$ root layer
(Lemma~\ref{lem:E8-to-H4-2to1}) transfers to the raw cumulative
phase $\Phiraw = 176\pi$ as an exact factor of $2$:
$\Phiraw \mapsto \Phiraw / 2$ under the $H_4$-observable phase
bookkeeping.
\end{hypothesis}

\begin{remark}[Status of $\Hphasehalve$]
\label{rmk:phase-halving-status}
The root-layer $2$-to-$1$ projection is a classical statement
about orbits of pairs of vectors $\{v, -v\}$; a transfer to the
\emph{cumulative} phase $\Phiraw$ is a statement about
integrating the sign-pairing over a plane-wise winding cycle
consistently with the $\sigma$-twist orientation $\Hor$ of
$\alpha$-2. Verifying this transfer requires a careful analysis
of how plane-wise windings at each Coxeter plane descend
through the sign-pair quotient, which we do not carry out
here. We adopt $\Hphasehalve$ as an explicit named hypothesis
of $\alpha$-3 (distinct from $\Hor, \Hdist, \Hquot, \Hfifty,
\Hgauge$).
\end{remark}

\begin{definition}[$H_4$-observable phase per cycle, conditional on $\Hphasehalve$]
\label{def:phi-H4}
Assuming Hypothesis~\ref{hyp:phase-halving}, the
\emph{$H_4$-observable phase per cycle} is
\[
  \PhiHfour \;:=\; \Phiraw / 2 \;=\; 176\pi / 2 \;=\; 88\pi.
\]
\end{definition}

\begin{lemma}[$c^{15} = -I$ in $W(E_8)$]
\label{lem:c15}
Let $c \in W(E_8)$ be a Coxeter element of order $h = 30$.
Then $c^{15} = -I$ acts as $-\mathrm{id}$ on the reflection
representation $V = \R^8$.
\end{lemma}

\begin{proof}
$c$ has order exactly $h = 30$ on $V$
\cite[\S 3.18]{Humphreys}. Its eigenvalues on $V \otimes \C$
are $\zeta^{m_k^\pm}$ with $\zeta = e^{2\pi i / 30}$ and
$\{m_k^\pm\} = \{1, 7, 11, 13, 17, 19, 23, 29\}$
(\cite[Lem.~lem:E8-data]{CascadePhasonCoxeter}). Then
$c^{15}$ has eigenvalues $\zeta^{15 m_k^\pm} = e^{i \pi
m_k^\pm}$. Since each exponent $m_k^\pm$ is odd (every
member of $\{1, 7, 11, 13, 17, 19, 23, 29\}$ is odd),
$e^{i \pi m_k^\pm} = -1$ for all $k$. Hence $c^{15}$ has
all eigenvalues equal to $-1$ on $V \otimes \C$, so
$c^{15} = -I$ on $V$.
\end{proof}

\begin{definition}[$\Z/2$ center generator]
\label{def:center}
The \emph{$\Z/2$ center generator} of $W(E_8)$ acting on
$V$ is the involution $z := c^{15} = -I$ of
Lemma~\ref{lem:c15}. The corresponding center quotient is
$V / \Z/2 = V / \langle -I \rangle$, which collapses
antipodal pairs.
\end{definition}

\begin{remark}[On the center quotient]
\label{rmk:center-pi}
A full Coxeter cycle on $V$ accumulates phase $2\pi$ in each
$2$D Coxeter plane $P_k$ (since $c|_{P_k}$ is rotation by
$2\pi m_k^- / 30$, and after $30$ steps the phase is
$2\pi m_k^-$, a multiple of $2\pi$). The center generator
$c^{15} = -I$ acts in each $P_k$ as a half-cycle: rotation
by $\pi$. Hence on each plane, the center generator
``accounts for'' a phase contribution of $\pi$ per cycle,
in the precise sense that one full Coxeter cycle equals
two applications of $c^{15}$ (since $c^{30} = (c^{15})^2 =
(-I)^2 = I$). At the level of accumulated phase
\emph{per slot}, this central involution is the natural
candidate for separating a $\pi$ contribution from the
fractional remainder; this candidacy is formalised as
$\Hgauge$ below.
\end{remark}

%=====================================================================
\section{Hypothesis $\Hgauge$: the $\pi$ center-gauge
  separation}
\label{sec:Hgauge}
%=====================================================================

\begin{hypothesis}[$\Hgauge$: $\pi$ separates as the
  $\Z/2$ center gauge]
\label{hyp:Hgauge}
Let $\PhiHfour = 88\pi$ be the $H_4$-observable phase per
cycle (Def.~\ref{def:phi-H4}) and $\Nslot = 87$ be the
phason slot count (Def.~\ref{def:Nslot}). The per-slot
phase
\[
  \PhiHfour / \Nslot \;=\; 88\pi / 87
\]
admits the canonical decomposition
\[
  88\pi / 87 \;=\; \pi + \pi/87,
\]
with the $+\pi$ term identified as the $\Z/2$ center-gauge
contribution generated by $c^{15} = -I$
(Lemma~\ref{lem:c15}, Remark~\ref{rmk:center-pi}). The
\emph{residual phase per slot} is
\[
  \Phi_{\rm res} \;:=\; (\PhiHfour / \Nslot) - \pi
                 \;=\; \pi/87.
\]
\end{hypothesis}

\begin{remark}[Status of $\Hgauge$]
\label{rmk:Hgauge-status}
$\Hgauge$ is a \emph{named structural hypothesis}, not a
theorem of $\alpha$-3. The arithmetic identity
$88\pi/87 = \pi + \pi/87$ is a pure algebraic computation
and requires no hypothesis. What $\Hgauge$ asserts beyond
the arithmetic is the \emph{identification} of the $+\pi$
piece with the $\Z/2$ center-gauge contribution generated
by $c^{15} = -I$.

The motivation
(Remark~\ref{rmk:center-pi}) is that the central involution
$c^{15}$ acts as a half-cycle (rotation by $\pi$) on each
Coxeter plane and naturally accounts for one half-revolution
per cycle. To upgrade $\Hgauge$ to a theorem one would
construct an explicit gauge-theoretic quotient of the
phason bookkeeping by the $\Z/2$ central action and prove
that this quotient removes exactly $\pi$ of phase per slot
(not, say, $0$ or $2\pi$). Such a construction would require
formalising the cut-and-project's phason bundle as a
$\Z/2$-equivariant object over the Coxeter base; this is
beyond the scope of $\alpha$-3.

In the meantime, $\Hgauge$ stays an explicit hypothesis
whose content is exactly the canonical separation of the
$\pi$ piece from the fractional remainder $\pi/87$ via the
center-gauge identification.
\end{remark}

\begin{theorem}[$\mathbf{T_{\pi/87}}$: per-slot residual
  phase; conditional on $\Hmin + \Hor + \Hdist + \Hquot +
  \Hphasehalve + \Hgauge$]
\label{thm:Tpi}
Assume the full $\alpha$-2 stack $\Hmin + \Hor + \Hdist +
\Hquot$, together with $\Hphasehalve$ (so that
Def.~\ref{def:phi-H4} gives $\PhiHfour = 88\pi$) and
$\Hgauge$. Then the per-slot $H_4$-observable phase decomposes
canonically as
\[
  \PhiHfour / \Nslot \;=\; 88\pi / 87 \;=\; \pi + \pi/87,
\]
with the $+\pi$ piece identified as the $\Z/2$ center
gauge (under $\Hgauge$); the residual phase per slot is
\[
  \Phi_{\rm res} \;=\; \pi/87.
\]
\end{theorem}

\begin{proof}
The identity $88\pi/87 = \pi + \pi/87$ is direct
arithmetic: $88\pi/87 = (87 + 1)\pi / 87 = \pi + \pi/87$.
The values $\PhiHfour = 88\pi$ and $\Nslot = 87$ are
inherited from the $\alpha$-2 stack via
Defs.~\ref{def:phi-H4}, \ref{def:Nslot} (carrying the
$\Hmin + \Hor + \Hdist + \Hquot$ fence). The
identification of the $+\pi$ piece with the $\Z/2$ center
gauge is the content of $\Hgauge$.
\end{proof}

%=====================================================================
\section{The Final Synthesis: Theorem $T_\alpha$}
\label{sec:Talpha}
%=====================================================================

\begin{definition}[$\alpha^{-1}$ as a derived constant, conditional on the full hypothesis stack]
\label{def:alpha-inv}
The symbol $\alpha^{-1}$ is a \emph{formal assembly} of three
ingredients of the $\alpha$-chain, each of which is available
only under the hypothesis stack named after it:
\[
  \alpha^{-1} \;:=\; (\Nslot + \kappaTwoI) + \Phi_{\rm res},
\]
where $\Nslot$ (the inherited phason slot count of
Def.~\ref{def:Nslot}) requires $\Hmin + \Hor + \Hdist + \Hquot$
(inherited from $\alpha$-2), $\kappaTwoI$ (the rational-sector
Laplacian sum of Def.~\ref{def:kappa}) requires $\Hfifty$, and
$\Phi_{\rm res}$ (the residual phase per slot) requires
$\Hphasehalve + \Hgauge$. Without these hypotheses, the value
of $\alpha^{-1}$ is not defined by this paper. \emph{Inside the
proof layer of $\alpha$-3, $\alpha^{-1}$ is a mathematical
symbol whose evaluated value is only available conditional on
the seven-hypothesis stack of Theorem~\ref{thm:Talpha}; no claim
about identification with the experimental fine-structure
constant is made within the theorem statements.} The identity
$\alpha^{-1} = 137 + \pi/87$ is not part of this definition; it
is the content of Theorem~\ref{thm:Talpha} under the
appropriate hypothesis stack.
\end{definition}

\begin{theorem}[$\mathbf{T_\alpha}$: the fine-structure
  identity, fully conditional]
\label{thm:Talpha}
Assume the seven-hypothesis stack
\[
  \Hmin \,+\, \Hor \,+\, \Hdist \,+\, \Hquot \,+\,
  \Hfifty \,+\, \Hphasehalve \,+\, \Hgauge.
\]
Then the derived constant $\alpha^{-1}$ of
Def.~\ref{def:alpha-inv} satisfies
\[
  \alpha^{-1} \;=\; 137 + \pi/87.
\]
\end{theorem}

\begin{proof}
By Theorem~\ref{thm:T137} (under $\Hmin + \Hor + \Hdist +
\Hquot + \Hfifty$), the integer part is $\Nslot +
\kappaTwoI = 87 + 50 = 137$. By Theorem~\ref{thm:Tpi}
(under $\Hmin + \Hor + \Hdist + \Hquot + \Hphasehalve +
\Hgauge$), the residual phase part is $\Phi_{\rm res} =
\pi/87$. The sum $137 + \pi/87$ is the value of $\alpha^{-1}$
by Def.~\ref{def:alpha-inv}. The full hypothesis stack is the
union $\Hmin + \Hor + \Hdist + \Hquot + \Hfifty + \Hphasehalve
+ \Hgauge$ as claimed.
\end{proof}

\begin{remark}[Status of $T_\alpha$]
\label{rmk:Talpha-status}
Theorem~\ref{thm:Talpha} is the main deliverable of
$\alpha$-3. It is conditional on seven named hypotheses;
none of them is itself a theorem of $\alpha$-3.
\begin{itemize}
\item $\Hmin$, $\Hor$, $\Hdist$, $\Hquot$ are inherited
  from $\alpha$-1 and $\alpha$-2 (where they are also
  named hypotheses, not theorems).
\item $\Hfifty$, $\Hphasehalve$, $\Hgauge$ are introduced
  for the first time in $\alpha$-3, with structural
  justifications in Remarks~\ref{rmk:H50-status},
  \ref{rmk:phase-halving-status}, \ref{rmk:Hgauge-status}.
\end{itemize}
What $\alpha$-3 \emph{does} prove is the conditional
arithmetic-and-phase identity that, given all seven
hypotheses, $\alpha^{-1} = 137 + \pi/87$. We make no claim
that this $\alpha^{-1}$ equals the experimental
fine-structure constant; that identification is outside the
proof layer of $\alpha$-3.
\end{remark}

%=====================================================================
\section{Self-Audit}
\label{sec:audit}
%=====================================================================

\textbf{Inherited (cited from $\alpha$-1 / $\alpha$-2;
not claimed as theorems of $\alpha$-3).}
\begin{itemize}
\item[[1]] $\alpha$-1 $C_1, C_2, C_3$
  (Fact~\ref{fact:imports}(a), (b1), (b2)) cited at their
  respective conditional levels ($C_2$ unconditional;
  $C_1, C_3$ on $\Hmin$).
\item[[2]] $\alpha$-2 $T_A$, $T_B$-$88$
  (Fact~\ref{fact:imports}(c)--(d)) cited conditionally on
  $\Hmin + \Hor + \Hdist$.
\item[[3]] $\alpha$-2 $T_B$-$87$
  (Fact~\ref{fact:imports}(e)) cited conditionally on
  $\Hmin + \Hor + \Hdist + \Hquot$.
\end{itemize}

\textbf{Classical layer (unconditional).}
\begin{itemize}
\item[[4]] Lemma~\ref{lem:McKay-recap} (McKay $2I
  \leftrightarrow \widetilde{E}_8$): McKay 1980 \S 3.
\item[[5]] Lemma~\ref{lem:c15} ($c^{15} = -I$): direct
  computation from the $E_8$ exponent table; all exponents
  $\{1, 7, 11, 13, 17, 19, 23, 29\}$ are odd, so $\zeta^{15
  m} = -1$ for $\zeta = e^{2\pi i/30}$.
\item[[6]] Lemma~\ref{lem:E8-to-H4-2to1} ($2$-to-$1$
  $E_8 \to H_4$): Moody--Patera Thm.~4.1, Prop.~4.2.
\item[[7]] Lemma~\ref{lem:2I-laplacian} ($2I$ rational
  irreps + Laplacian formula): Du Val Ch.~3 (character table)
  + in-paper Proposition~\ref{prop:class-sum} (class-sum /
  Cayley-graph spectral formula, proved here).
\end{itemize}

\textbf{New named structural hypotheses (NOT theorems of
$\alpha$-3).}
\begin{itemize}
\item[[8]] $\Hfifty$ (Hypothesis~\ref{hyp:H50}):
  $\kappaTwoI = 50$; structural motivation in
  Remark~\ref{rmk:H50-status}, with the gap explicitly
  identified (convention pinning + finite computation).
\item[[9]] $\Hgauge$ (Hypothesis~\ref{hyp:Hgauge}): the
  $+\pi$ piece of $88\pi/87$ is canonically the $\Z/2$
  center-gauge contribution; structural motivation in
  Remarks~\ref{rmk:center-pi}, \ref{rmk:Hgauge-status},
  with the gap explicitly identified
  ($\Z/2$-equivariant phason-bundle quotient construction).
\end{itemize}

\textbf{Conditional theorems of $\alpha$-3.}
\begin{itemize}
\item[[10]] Theorem~\ref{thm:T137} ($T_{137}$, integer
  decomposition $137 = 87 + 50$) conditional on
  $\Hmin + \Hor + \Hdist + \Hquot + \Hfifty$. Substantive
  content lives in the hypothesis stack
  (Remark~\ref{rmk:T137-status}).
\item[[11]] Theorem~\ref{thm:Tpi} ($T_{\pi/87}$, per-slot
  residual phase) conditional on $\Hmin + \Hor + \Hdist +
  \Hquot + \Hphasehalve + \Hgauge$. Arithmetic identity is
  direct; substantive content is the phase-halving transfer
  ($\Hphasehalve$) and center-gauge identification
  ($\Hgauge$).
\item[[12]] Theorem~\ref{thm:Talpha} ($T_\alpha$,
  $\alpha^{-1} = 137 + \pi/87$) conditional on the full
  seven-hypothesis stack
  $\Hmin + \Hor + \Hdist + \Hquot + \Hfifty + \Hphasehalve + \Hgauge$.
\end{itemize}

\textbf{Not claimed.}
\begin{itemize}
\item No unconditional $T_\alpha$.
\item No claim that $\alpha^{-1}$ as defined equals an
  experimental observable.
\item No phenomenology, photons, microtubules, biology,
  Standard Model, or first-principles physics rhetoric.
\item No CODATA / measurement comparison inside the proof
  layer.
\item No dependence on P1--P7, on \texttt{cascade-derivation/}
  working notes, on \texttt{paper-xxii},
  \texttt{paper-xxxiv}, or any other downstream VFD
  manuscript.
\item No GAP / computer-algebra session is cited as proof
  source.
\end{itemize}

%=====================================================================
\section{Citation Contract and Downstream Handoff}
\label{sec:handoff}
%=====================================================================

\textbf{Unconditionally citable.}
\begin{itemize}
\item Lemmas~\ref{lem:McKay-recap}, \ref{lem:c15},
  \ref{lem:E8-to-H4-2to1}, \ref{lem:2I-laplacian}
  (classical layer).
\end{itemize}

\textbf{Conditional on $\alpha$-1's $\Hmin$.}
\begin{itemize}
\item $\alpha$-1's $C_1$, $C_3$ (cited as inputs).
\end{itemize}

\textbf{Conditional on $\Hmin + \Hor + \Hdist$.}
\begin{itemize}
\item $\alpha$-2's $T_A$, $T_B$-$88$ (cited as inputs);
  $\Phiraw$, $\PhiHfour$ definitions.
\end{itemize}

\textbf{Conditional on $\Hmin + \Hor + \Hdist + \Hquot$.}
\begin{itemize}
\item $\alpha$-2's $T_B$-$87$ (cited as input); $\Nslot$
  definition.
\end{itemize}

\textbf{Conditional on $\Hmin + \Hor + \Hdist + \Hquot +
\Hfifty$.}
\begin{itemize}
\item Theorem~\ref{thm:T137}: $\Nslot + \kappaTwoI = 137$.
\end{itemize}

\textbf{Conditional on $\Hmin + \Hor + \Hdist + \Hquot +
\Hphasehalve + \Hgauge$.}
\begin{itemize}
\item Theorem~\ref{thm:Tpi}: per-slot residual $\pi/87$.
\end{itemize}

\textbf{Conditional on the full seven-hypothesis stack.}
\begin{itemize}
\item Theorem~\ref{thm:Talpha}: $\alpha^{-1} = 137 + \pi/87$.
\end{itemize}

\textbf{Does NOT close.}
\begin{itemize}
\item Any of the seven hypotheses as a theorem of $\alpha$-3.
  $\Hmin$, $\Hor$, $\Hdist$, $\Hquot$ remain open from
  $\alpha$-1 / $\alpha$-2; $\Hfifty$, $\Hphasehalve$, and
  $\Hgauge$ are introduced here as named hypotheses with
  structural motivation but not theorem-grade derivations.
\item Any identification of $\alpha^{-1}$ with an
  experimental observable.
\end{itemize}

\textbf{Downstream usage.} Any future paper wishing to
identify the derived constant $\alpha^{-1} = 137 + \pi/87$
of $\alpha$-3 with the experimentally measured fine-structure
constant must:
\begin{enumerate}
\item Either upgrade some or all of the seven hypotheses to
  theorems, or assume them explicitly with downstream
  conditional fences.
\item Bridge the mathematical $\alpha^{-1}$ of
  Def.~\ref{def:alpha-inv} to a physical observable, which
  is outside the strictly mathematical scope of $\alpha$-3.
\end{enumerate}

%=====================================================================
\section{Notation Audit}
\label{sec:notation}
%=====================================================================

\begin{center}
\begin{tabular}{|l|l|l|}
\hline
Symbol & Meaning & First use \\
\hline
$\Nslot$ & phason slot count $= 87$ (from $\alpha$-2) &
  Def.~\ref{def:Nslot} \\
$2I$ & binary icosahedral group &
  Lemma~\ref{lem:McKay-recap} \\
$\Irrrat(2I)$ & $2I$ irreducibles with rational character &
  Def.~\ref{def:Irr-rat} \\
$\Irrrat^\times(2I)$ & nontrivial part of $\Irrrat(2I)$ &
  Def.~\ref{def:Irr-rat} \\
$L = 12 I - A$ & $600$-cell graph Laplacian &
  Def.~\ref{def:Lambda} \\
$\lambda_\rho$ & eigenvalue of $L$ on $V_\rho$-isotypic
  component & Def.~\ref{def:Lambda} \\
$\kappaTwoI$ & $\sum_{\rho \in \Irrrat^\times(2I)}
  \lambda_\rho$ & Def.~\ref{def:kappa} \\
$\Wph$ & total phason winding $= 88$ (from $\alpha$-2) &
  Fact~\ref{fact:imports}(d) \\
$\Phiraw$ & $\Wph \cdot 2\pi = 176\pi$ &
  Def.~\ref{def:phi-raw} \\
$\PhiHfour$ & $\Phiraw / 2 = 88\pi$, $H_4$-observable phase
  per cycle & Def.~\ref{def:phi-H4} \\
$c$ & Coxeter element of $W(E_8)$, order $30$ &
  Lemma~\ref{lem:c15} \\
$z = c^{15}$ & $\Z/2$ center generator, $-I$ on $V$ &
  Def.~\ref{def:center} \\
$\Phi_{\rm res}$ & residual phase per slot $= \pi/87$ under
  $\Hgauge$ & Hyp.~\ref{hyp:Hgauge} \\
$\alpha^{-1}$ & derived constant $= 137 + \pi/87$ under
  full stack & Def.~\ref{def:alpha-inv} \\
$\Hmin, \Hor, \Hdist, \Hquot$ & inherited hypotheses
  ($\alpha$-1, $\alpha$-2) &
  Fact~\ref{fact:imports} \\
$\Hfifty$ & $\kappaTwoI = 50$ identification &
  Hyp.~\ref{hyp:H50} \\
$\Hgauge$ & $+\pi$ separates as $\Z/2$ center gauge &
  Hyp.~\ref{hyp:Hgauge} \\
$T_{137}, T_{\pi/87}, T_\alpha$ &
  Theorems~\ref{thm:T137}, \ref{thm:Tpi},
  \ref{thm:Talpha} & \\
\hline
\end{tabular}
\end{center}

\begin{thebibliography}{99}

\bibitem{CascadeFoundations}
\emph{Cascade Correspondence Foundations},
\texttt{papers/cascade-correspondence-foundations/foundations.tex}.
Cited solely for the statement of Axiom $\mathbf{F1}$ (\S F1).

\bibitem{Cascade12D}
\emph{The 12D Closure Theorem} ($\alpha$-1),
\texttt{papers/cascade-12d-closure/cascade-12d-closure.tex}.
Theorems $C_1$, $C_2$, $C_3$ and Hypothesis $\Hmin$.

\bibitem{CascadePhasonCoxeter}
\emph{Phason--Coxeter Alignment} ($\alpha$-2),
\texttt{papers/cascade-phason-coxeter/cascade-phason-coxeter.tex}.
Theorems $T_A$, $T_B$-$88$, $T_B$-$87$ and Hypotheses
$\Hor$, $\Hdist$, $\Hquot$.

\bibitem{Bourbaki46}
N.~Bourbaki, \emph{Groupes et alg\`ebres de Lie}, Chapitres
4--6, Hermann/Springer, 1968. $E_8$ data: Ch.~6, \S4 No.~12,
Planche~VII; Coxeter elements and exponent symmetry: Ch.~5,
\S6.

\bibitem{Humphreys}
J.~E.~Humphreys, \emph{Reflection Groups and Coxeter Groups},
Cambridge Stud. Adv. Math.~29, Cambridge UP, 1990.
\S 3.16 (Coxeter element); \S 3.18 (Coxeter number);
\S 3.19--3.20 (eigenvalues / exponents).

\bibitem{ConwaySloane}
J.~H.~Conway, N.~J.~A.~Sloane, \emph{Sphere Packings,
Lattices and Groups}, 3rd ed., Springer, 1999. Icosian
ring and $E_8$ embedding: Ch.~8, \S2 (pp.~206--209);
$600$-cell: Ch.~8, \S2.3.

\bibitem{MoodyPatera}
R.~V.~Moody, J.~Patera, \emph{Quasicrystals and icosians},
J.~Phys.~A \textbf{26} (1993), 2829--2853. Theorem~4.1 and
Proposition~4.2 (the $E_8 \to H_4$ cut-and-project, including
the $2$-to-$1$ observable structure).

\bibitem{McKay80}
J.~McKay, \emph{Graphs, singularities, and finite groups},
in: \emph{Finite Groups (Santa Cruz, 1979)},
Proc.~Symp.~Pure Math.~37, AMS, 1980, pp.~183--186.

\bibitem{DuVal}
P.~Du~Val, \emph{Homographies, Quaternions and Rotations},
Oxford Mathematical Monographs, Clarendon Press, 1964.
Ch.~3: $2I$ and its character table.

\end{thebibliography}

\end{document}
